\section{Representaciones abstractas: razonamiento médico y circuitos multilingües}

Una vez establecido el método, el primer hallazgo es que dentro del modelo existen representaciones abstractas que se reutilizan entre distintas formas superficiales. Esta sección lo ilustra con dos casos, uno médico y otro multilingüe: un mismo concepto puede activarse con palabras distintas, y la información de frontera sobre el idioma de entrada y el de salida puede separarse de la parte semántica intermedia.

\subsection{Caso médico: síntomas, diagnóstico y criterios se combinan en una sola propagación hacia adelante}

El caso clínico presenta embarazo en el tercer trimestre, dolor en el cuadrante superior derecho del abdomen, hipertensión y enzimas hepáticas elevadas, y pide preguntar por un solo síntoma adicional. El modelo responde visual disturbances. Lo importante no es la conclusión médica en sí, sino si en el grafo aparecen representaciones combinables como pregnancy, preeclampsia y diagnostic criteria, y si la respuesta cambia de la manera esperada al suprimir alguna feature.

\lecturefigure{slide-25-medical-reasoning-prompt.jpg}{Prompt de razonamiento médico en una sola propagación hacia adelante}{00:37:30}

\lecturefigure{slide-26-medical-concept-representations.jpg}{El modelo organiza los síntomas, el diagnóstico y los criterios diagnósticos adicionales como features abstractas}{00:38:46}

En la figura 26, varias features de síntomas convergen primero en una representación relacionada con la enfermedad, y después la representación del diagnóstico impulsa «se debe preguntar por visual disturbances». No se trata de una chain of thought en lenguaje natural, pero sí refleja relaciones estructuradas: los síntomas no se emparejan palabra por palabra, sino que confluyen internamente como evidencia diagnóstica.

\lecturefigure{slide-27-medical-feature-interventions.jpg}{Al sustituir o inhibir features médicas, la pregunta de seguimiento recomendada cambia en consecuencia}{00:39:51}

La intervención aumenta la credibilidad causal: si al elevar las features relacionadas con el biliary-system la salida se desplaza hacia otro tipo de síntoma, esas direcciones no son solo nombres asignados a posteriori. Sin embargo, el conocimiento clínico puede estar incompleto y el prompt es muy simplificado; la figura no demuestra que el modelo sea apto para el diagnóstico real, y mucho menos puede sustituir a las guías médicas ni al juicio del médico.

\begin{warningbox}{Alcance del caso médico}
Este es un mechanistic case study, no un consejo médico ni una prueba de confiabilidad clínica. Muestra que el modelo puede combinar representaciones médicas en un prompt; no muestra que su conocimiento sea completo, que sus probabilidades estén calibradas, que sea equitativo entre poblaciones ni que sea seguro en escenarios reales.
\end{warningbox}

\teachervoice{00:31:52--00:42:44: el docente usa el ejemplo médico como una ventana a «cómo un solo forward pass realiza razonamiento composicional», y en la sesión de preguntas reconoce que la confiabilidad en el ámbito médico real requiere una verificación independiente.}

\subsection{Caso multilingüe: en la frontera está el idioma, en el núcleo la semántica compartida}

«El antónimo de small» puede preguntarse en inglés, francés o chino. Si el modelo pensara por completo en inglés, todo el grafo interno tendría que traducir primero al inglés y después calcular; si cada idioma fuera totalmente independiente, los tres grafos casi no compartirían nada. El resultado real se parece más a una estructura por capas: cerca de la entrada y de la salida hay features específicas de cada idioma, y en la parte intermedia aparecen representaciones compartidas como multilingual small, antonym y large.

\lecturefigure{slide-28-unique-language-question.jpg}{Pregunta: ¿el modelo «piensa» en un único idioma?}{00:42:47}

\lecturefigure{slide-29-shared-multilingual-representations.jpg}{Los tres idiomas comparten features semánticas intermedias y conservan fronteras específicas de cada idioma}{00:45:01}

Al leer la figura, primero se miran los dos extremos: quote, English/Chinese/French y los tokens concretos se encargan de la forma lingüística; después se mira la parte intermedia: features como small, opposite y say large coinciden entre idiomas. La conclusión no es que «el modelo tiene un lenguaje puro del pensamiento», sino que la abstracción semántica y la realización superficial están parcialmente desacopladas.

\lecturefigure{slide-30-antonym-to-synonym-edit.jpg}{Al sustituir la feature antonym por la feature synonym, la salida en los tres idiomas cambia según la semántica}{00:45:06}

Esta intervention es decisiva: una misma edición interna convierte, en todos los idiomas, la tarea de «antónimo» en la de «sinónimo», lo que indica que las features compartidas participan en el cálculo y no solo están correlacionadas con la salida. El operand de entrada, el tipo de operación y el idioma de salida pueden editarse por separado, lo que revela una representación interna modular, aunque no completamente discreta.

Si los conjuntos de features activas en dos idiomas son $F_a,F_b$, la proporción compartida puede expresarse con la intersección sobre la unión:
\[
\operatorname{IoU}(F_a,F_b)=\frac{|F_a\cap F_b|}{|F_a\cup F_b|}.
\]
La intersección cuenta las features comunes y la unión cuenta las features que aparecen en al menos un idioma. Un IoU más alto indica más representaciones compartidas, pero depende del feature matching, del umbral y de la calidad del modelo de reemplazo.

\lecturefigure{slide-31-representation-sharing-scales.jpg}{Los modelos más grandes comparten una proporción mayor de representaciones entre idiomas}{00:47:26}

La tendencia de la figura 31 respalda que «el aumento de escala favorece la reutilización de abstracciones», pero no demuestra que todos los conceptos converjan en un mismo espacio semántico. La distancia entre idiomas, la tokenización, el volumen de datos y las expresiones culturalmente específicas siguen generando diferencias; la curva es el feature overlap en este diseño experimental, no una ley lingüística universal.

\begin{knowledgebox}{De «el idioma del pensamiento» a tres preguntas medibles}
No conviene preguntar en qué idioma «piensa» realmente el modelo. Se pueden plantear por separado tres preguntas: en la etapa de entrada, qué features dependen de la forma lingüística; si las operaciones intermedias comparten direcciones semánticas; y cómo se elige el idioma de destino en la etapa de salida. Estas tres preguntas pueden verificarse con grafos e intervenciones, mientras que una sola pregunta basada en la atribución de rasgos humanos tiende a generar confusión.
\end{knowledgebox}

\subsection{Resumen de la sección}

Los casos médico y multilingüe muestran en conjunto que el modelo combina tokens superficiales en conceptos y operaciones más abstractos. Las representaciones abstractas pueden reutilizarse entre síntomas, idiomas y expresiones, pero siguen inscritas en un prompt, un modelo y un método de extracción de features concretos; que sean «reutilizables» no permite concluir que sean «universales, únicas y completas».
